%% appB-notation \dash{} Notacja i symbole

\chapter{Notacja i symbole}
\label{app:B}

\noindent
Ten dodatek odnotowuje cztery miejsca, w których uzus matematyczny jest
podzielony i książka musiała wybrać, oraz jedną regułę rozstrzygającą, jak
nazwa jest zapisywana w każdym z wydań. Celowo nie jest tabelą wszystkich
symboli: czytelnik, który nigdy nie widział $\sum$, spotyka go zbudowanego od
zera w Programie~\ref{prog:F04}, a indeks pod nim w Programie~\ref{prog:F02},
więc robotę, którą miałaby wykonać tabela symboli, wykonują programy
wprowadzające te symbole.

% ============================================================
\section{Decyzje, które ta książka podjęła}
\label{sec:appB-decisions}

\noindent
Tam, gdzie uzus matematyczny jest podzielony, książka musi wybrać, a sam wybór
jest dla czytelnika wart więcej niż konwencja. Oto cztery, które naprawdę
bolą.

\begin{notationbox}
\textbf{Logarytm niesie tutaj swoją podstawę, a build tego pilnuje.}
Logarytm naturalny to $\ln$, logarytm o podstawie dwa to $\logb{2}$. Sam
$\mfalogplain$ jest błędem builda, a lint parzystości nazywa plik i wiersz,
zanim build w ogóle do tego dojdzie.

Są dwa rodzaje miejsc, w których zakaz jest uchylony, i oba są celowe. Jedno
to miejsce, w którym sam zakaz jest tematem: ten dodatek oraz sekcja~2
Programu~\ref{prog:F03}, która każe czytelnikowi spotkać to zderzenie, zanim
zostanie nazwane. Drugie to \textbf{wnętrze $O$}, gdzie zmiana podstawy mnoży
przez stałą, a stała jest właśnie tym, co ta notacja odrzuca \dash{} więc
$O(n \mfalogplain n)$ nazywa jedną klasę wzrostu naraz dla każdej podstawy, a
wpisanie tam podstawy sugerowałoby rozróżnienie, którego nie ma.
Program~\ref{prog:P03} ma to jako ramkę notacyjną, a Program~\ref{prog:P16} z
tego korzysta. Liczba takich miejsc celowo nie jest tu podana, bo takie
zliczenie się starzeje, a nic nie może go sprawdzić.

Powód jest taki, że dwie grupy czytelników, którym ta książka służy, zostały
nauczone przeciwnych konwencji. W polskich podręcznikach sam logarytm oznacza
podstawę dziesięć. W piśmiennictwie uczenia maszynowego oznacza podstawę $e$.
Te same trzy znaki niosą przeciwne znaczenia i zderzają się najmocniej
dokładnie w tych programach, w których najbardziej się to liczy: entropię
podaje się w bitach, a entropię krzyżową w natach, a dzielą je dwa programy.

Konwencja ogłoszona raz w dodatku tego nie naprawia. Napisanie podstawy
kosztuje dwa znaki, a build odrzuca sam logarytm, zamiast liczyć na czyjąś
pamięć.
\end{notationbox}

\begin{notationbox}
\textbf{Wektory są kolumnami.} $x \in \R^{n}$ jest kolumną, a $x\T$ wierszem.
Polska szkoła pisze też strzałkę nad literą; prace, do których czytania ta
książka cię przygotowuje, tego nie robią, więc ta książka też nie.

Wektor składa się zwykłą kursywą wszędzie tam, gdzie nic innego w wyrażeniu
nie mogłoby zostać za niego wzięte, czyli w większości książki \dash{}
Programy \ref{prog:F09}, \ref{prog:P04} i \ref{prog:P05} wprowadzają wektory i
składają każdy z nich kursywą. \textbf{Pogrubienia} używamy tam, gdzie wektor
i macierz dzielą jedno wyrażenie, a chodzi właśnie o kształty, jak w
$\vect{y} = W\vect{x}$: to jest temat Programu~\ref{prog:P18} od początku do
końca i tych kilku późniejszych miejsc, które go cytują.

\textbf{Pochodne macierzowe używają układu licznikowego}, zadeklarowanego w
Programie~P18 i nigdy nie przyjmowanego milcząco. To nie jest różnica
polsko-angielska, tylko różnica między dziedzinami, i bierze się z niej dobra
połowa zamieszania z transpozycjami w opisach propagacji wstecznej.
\end{notationbox}

\begin{notationbox}
\textbf{Przedziały zapisujemy nawiasami kwadratowymi:} $\intcc{a}{b}$ jest
domknięty, $\intoo{a}{b}$ otwarty. Polska szkoła zapisuje przedział domknięty
jako $\langle a, b\rangle$. Norma ISO 80000-2 i każda praca, którą otworzysz,
tego nie robią. Książka idzie za pracami. Pierwszy przedział, jaki czytelnik
spotyka, jest w Programie~F3, a żaden program nie zatrzymuje się, żeby ten
wybór wyjaśnić, więc odnotowuje go ten dodatek.
\end{notationbox}

\begin{notationbox}
\textbf{Wariancja to $\Var(X)$, a wartość oczekiwana to $\Ex[X]$.} Polskie
materiały dydaktyczne do dziś piszą $D^{2}(X)$ obok $\Var$ \dash{} to jest bieżący
uzus, a nie ciekawostka historyczna, i Program~P24 podaje tę równoważność tam,
gdzie pojęcie się pojawia. $M(X)$ na wartość oczekiwaną jest formą reliktową i
dostaje zdanie, a nie ramkę.
\end{notationbox}

% ============================================================
\section{Test klawiatury}
\label{sec:appB-keyboard}

\noindent
Wydanie polskie pisze \emph{tg} na tangens, tam gdzie wydanie angielskie
pisze \emph{tan} \dash{} a oba wydania piszą \code{tanh} na tangens
hiperboliczny, zamiast \emph{tgh}, które składa polskie piśmiennictwo
matematyczne. Wygląda to na niekonsekwencję. To jest reguła w działaniu.

Obie pisownie są tu złożone jako dosłowny tekst, a nie przez makra notacyjne,
i musi tak być: makro renderuje się w wydaniu, które trzymasz, więc napisanie
tego akapitu przez $\tg$ dałoby w obu kolumnach to samo słowo i opisywałoby
rozbieżność, której strona nie pokazuje.

\begin{note}
\textbf{Jeśli będziesz wpisywał dany zapis do komputera, oba wydania używają
pisowni, którą wpiszesz. Jeśli będziesz go wyłącznie czytał albo pisał ręcznie,
wydanie polskie używa formy polskiej.}

Nikt nie wpisuje \emph{tg} do programu, więc wydanie polskie składa
\emph{tg}, \emph{ctg} i \emph{arc tg} tak, jak robi to polska matematyka
\dash{} przez $\tg$, $\ctg$ i $\arctg$, które tutaj składają formy polskie, a
tam angielskie. Ale \code{tanh} to
\code{torch.tanh} \dash{} jest funkcją aktywacji, zanim jest funkcją zmiennej
rzeczywistej, a polski czytelnik, który nauczy się \emph{tgh}, nie znajdzie jej
w żadnej bibliotece. Polskie piśmiennictwo matematyczne pisze \emph{tgh}, a
żaden program nie zatrzymuje się, żeby to powiedzieć, co czyni ten dodatek
zapisem tej informacji.
\end{note}
